% \documentclass[polish, 8pt]{article}
\documentclass[polish, 10pt]{extarticle}
\usepackage[top=2cm, bottom=2cm, left=3cm, right=3cm]{geometry}
\usepackage{graphicx} % Required for inserting images

\usepackage[polish]{babel}
\usepackage{polski}
\usepackage{placeins}
\usepackage{dirtytalk}
% Set page size and margins
% Replace `letterpaper' with `a4paper' for UK/EU standard size
\usepackage{csquotes}
% Useful packages
\usepackage{float}
\usepackage{amsmath}
\usepackage{graphicx}
\usepackage{hyperref}
\usepackage{graphicx}
\usepackage{listings}
\usepackage{xcolor}
\usepackage{amssymb}
\usepackage{enumitem}
\usepackage{float}
\definecolor{codegreen}{rgb}{0,0.6,0}
\definecolor{codegray}{rgb}{0.5,0.5,0.5}
\definecolor{codepurple}{rgb}{0.58,0,0.82}
\definecolor{backcolour}{rgb}{0.95,0.95,0.92}

\lstdefinestyle{mystyle}{
    backgroundcolor=\color{backcolour},   
    commentstyle=\color{codegreen},
    keywordstyle=\color{magenta},
    numberstyle=\tiny\color{codegray},
    stringstyle=\color{codepurple},
    basicstyle=\ttfamily\footnotesize,
    breakatwhitespace=false,         
    breaklines=true,                 
    captionpos=b,                    
    keepspaces=true,                 
    numbers=left,                    
    numbersep=5pt,                  
    showspaces=false,                
    showstringspaces=false,
    showtabs=false,                  
    tabsize=2
}
\lstset{style=mystyle}
\hypersetup{
    colorlinks=true,
    linktoc=all     %set to all if you want both sections and subsections linked
}
\usepackage{mdframed}
\usepackage{titling}
\setlength{\droptitle}{-2cm} 

\begin{document}

\title{Kolokwium I -- Zestaw B}
\author{}
\date{}

\maketitle
\vspace{-2cm}

\section*{Zadanie 1 (25 pkt)}
Planujesz rozegrać dwie partie szachów z nieznanym wcześniej przeciwnikiem. Zakładamy, że Twój rywal reprezentuje jeden z trzech poziomów: początkujący, średniozaawansowany lub mistrzowski, przy czym każdy z tych poziomów jest jednakowo prawdopodobny.
Twoje szanse na wygraną w pojedynczej partii wynoszą odpowiednio: \(90\%\) (przeciwko początkującemu), \(50\%\) (przeciwko średniozaawansowanemu) oraz \(30\%\) (przeciwko mistrzowi).

\begin{enumerate}
  \item Oblicz prawdopodobieństwo wygrania pierwszej partii.
  \item Załóżmy, że pierwsza partia zakończyła się Twoją wygraną. Oblicz, jakie jest teraz prawdopodobieństwo, że Twój przeciwnik jest mistrzem. Jakie jest prawdopodobieństwo wygrania również drugiej partii, wiedząc że pierwsza zakończyła się wygraną? (Przyjmij założenie, że przy ustalonym poziomie umiejętności przeciwnika wyniki poszczególnych partii są od siebie niezależne). 
  \item Wyjaśnij różnicę między założeniem, że wyniki partii są \textit{niezależne}, a założeniem, że są one \textit{niezależne warunkowo} względem poziomu umiejętności przeciwnika. Które z tych założeń jest w tym kontekście bardziej uzasadnione i dlaczego?
\end{enumerate}

\section*{Zadanie 2 (15 pkt)}
Niech \(A\) i \(B\) będą zdarzeniami losowymi.
Różnicą symetryczną \(A \triangle B\) nazywamy zdarzenie zachodzące wtedy, gdy zachodzi \(A\) lub zachodzi \(B\), ale nie oba jednocześnie (w logice odpowiada to alternatywie wykluczającej \emph{XOR}). Korzystając z definicji i własności prawdopodobieństwa, wykaż, że:
\begin{enumerate}
\item $P(A \triangle B) \;=\; P(A) + P(B) - 2P(A \cap B)$
\item $P(A) - P(B) \leq P(A \triangle B)$
\end{enumerate}


\section*{Zadanie 3 (10 pkt)}
Podaj definicję dystrybuanty rozkładu zmiennej losowej rzeczywistej (jednowymiarowej) oraz
przynajmniej 4 z podanych na wykładzie podstawowych własności dystrybuant jednowymiarowych.

\section*{Zadanie 4 (30 pkt)}

Niech $X$ będzie zmienną losową o nieznanym parametrze $\theta > 0$, której funkcja
gęstości prawdopodobieństwa ma postać
\[
f_\theta(x) =
\begin{cases}
\dfrac{2x}{\theta^2}, & \text{dla } 0 \leq x \leq \theta,\\[4pt]
0, & \text{w przeciwnym przypadku.}
\end{cases}
\]

\begin{enumerate}[label=\alph*)]
    \item Dla jakich wartości parametru $\theta$ funkcja $f_\theta$ jest
    funkcją gęstości prawdopodobieństwa zmiennej losowej $X$?

    \item Wyznacz dystrybuantę $F_X(x)$ zmiennej losowej $X$.

    \item Oblicz wartość oczekiwaną $E[X]$, $Var(2X + 120)$ oraz medianę
    zmiennej losowej $X$ w zależności od parametru $\theta$.

    \item Oblicz prawdopodobieństwo $P\bigl(X^2 > 1\bigr)$ w zależności od parametru $\theta$.
\end{enumerate}

\section*{Zadanie 5 (20 pkt)}
Zmienna losowa \(X\) ma rozkład jednostajny \(U(-2, 3)\).
Wyznacz rozkład zmiennej losowej \(Y = (X-1)^2\).

\end{document}

